\section{A precise route to complete Lean formalization}
\label{sec:lean}

\subsection{Three separate obligations}

A semantic Lean theorem about \texttt{Nat}, lists, or automata does not by
itself prove that the theorem is derivable in iPRA. The completion plan
therefore has three separate obligations for every operation and theorem.

\begin{enumerate}
\item \textbf{Semantic mathematics.} Define the finite structures and
prove their ordinary Lean correctness statements. Finite alphabets and
state spaces can be \texttt{Fin k}; words are lists; finite relations
are Boolean tables. Quantified assertions about an arbitrary infinite
set of words are avoided in the numerical layer.
\item \textbf{Primitive-recursive representation.} Give numerical codecs,
well-formedness bits, and a fixed primitive-recursive expression for every
operation. Prove round trips for valid codes and prove that the coded
operation agrees with the semantic one. A function that Lean accepts as
total is not thereby known to be primitive recursive.
\item \textbf{Object-theory derivability.} Construct
\texttt{Prf PRAAx \(\Gamma\) \(\varphi\)} for the first-order formula
expressing the actual theorem. Every induction node must use a
quantifier-free formula. Bounded predicates are replaced by their bits
using the inspected \texttt{BoundedFormulas} infrastructure. The final
theorem has empty context and an explicit universally closed sentence.
\end{enumerate}

These tasks are individually substantial. The supplied files specify the
data, algorithms, invariants, witness bounds, and remaining proof targets;
they do not fill an axiom-shaped placeholder and count it as a theorem.
The report's proofs supply the mathematical content to be translated.

\subsection{Finite data and total conventions}

\begin{longtable}{@{}p{0.18\textwidth}p{0.37\textwidth}p{0.38\textwidth}@{}}
\toprule
Object & Semantic data & Numerical representation and check\\
\midrule\endhead
Alphabet and word & \texttt{List (Fin k)} & Base-$(k+1)$ word code,
positive digits, explicit length and validity bit.\\
DFA & Start, transition table, final bits & Positive state count, exactly
$nk$ table entries below $n$, exactly $n$ final bits.\\
$\eps$-NFA & Initial/final bit sets and tagged edges & Letter-or-empty
edge tags with checked endpoints; $2^n$ subset codes.\\
Regular expression & Inductive expression tree & Postorder node table;
children precede parents; each nonroot has one parent; all nodes reach
the root. Sharing must be unfolded before the linear Thompson bound.\\
Grammar & Tagged variables/terminals and rule list & Checked left-side
variables, right-side symbol bounds, start and finite rule indices.\\
Parse tree & Rule-labelled ordered rooted tree & Postorder table with
unique parents, valid local productions, and computed yields.\\
Derivation & Finite rewriting sequence & Forms, positions, and rule
indices; bounded adjacent-step checker.\\
PDA & Finite transition list, stack words & Input-or-empty tags,
pop symbol, pushed word, states, and complete-run certificate.\\
Machine & Finite table and configurations & State and two base-$b$ tape
halves; fixed Boolean output convention; bounded iteration.\\
Monoid & Finite multiplication table & Unit and associative table check;
letter map and accepting bit vector.\\
Polynomial & Finite signed coefficient list & Integer pairs and finite
convolution; degree and coefficient bounds.\\
\bottomrule
\end{longtable}

Malformed numerical descriptions receive explicit defaults. Theorems about
the intended semantics retain a validity premise. A rejecting default is
used for language descriptors; the machine chapter uses its stated
immediate-halt default for malformed machine codes. Those conventions
are independent and must not be conflated by an adapter.

\subsection{Representative endpoint formulas}

Here are the actual shapes of the closed arithmetic targets. Each symbol
in them denotes a fixed primitive-recursive function or a displayed
first-order formula; no extra sorts are implicit.

\paragraph{Determinization.}
\[
 \forall E,w\;\bigl(\operatorname{ENFAValid}(E)\land
 \operatorname{WordValid}(\operatorname{alphabet}(E),w)\bigr)
 \to \operatorname{Accept}_{\rm DFA}(\operatorname{Det}(E),w)
      =\operatorname{Accept}_{\eps\rm NFA}(E,w).
\]
Use the prefix invariant that each subset bit is exactly the corresponding
reachable-state bit. The invariant quantifies only over the supplied finite
state set and the processed prefix. Normalize empty-edge blocks to fewer
than $n$ edges before constructing or comparing path witnesses.

\paragraph{Equivalence of finite automata.}
For valid $A,B$ on the same alphabet, define a bit $\operatorname{Eq}(A,B)$
by the product reachability table. Prove
\[
 \operatorname{Eq}(A,B)=1\ \longleftrightarrow\
 \forall w\bigl(\operatorname{WordValid}(k,w)\to
                  \operatorname{Accept}(A,w)=\operatorname{Accept}(B,w)\bigr).
\]
The right-hand universal formula is not an induction motive. First prove
the finite fixed-point table properties. Then fix $w$ and induct through
its run with a bounded table-membership invariant. If the bit is zero,
the stored predecessor table supplies a separator of length less than
$|A||B|$. This supplies a concrete witness to inequivalence.

\paragraph{Grammar membership.}
For valid $G,w$, the target is
\[
 \operatorname{Parse}(G,w)=1\ \longleftrightarrow\
 \exists t\;\operatorname{Tree}(G,S,w,t)=1.
\]
The forward direction exhibits the parser's tree term. The reverse direction
fixes an arbitrary supplied tree and proves the table invariant by a scan
of its finite code. The CNF version strengthens the existential to a
primitive-recursively bounded one: the selected normal tree has $3|w|-1$
nodes for nonempty input and the explicit inverse-normalization bound
controls the original-grammar tree. Turn a bound on node count and labels
into a bound on its numerical code using the chosen table codec. Empty
input is a separate branch, never an instance of the formula $3|w|-1$.

\paragraph{Pumping.}
Compute the parser's chosen tree, a longest path, repeated-variable
positions, and the five cut positions. The first endpoint proves that
these cuts satisfy the bounded conditions. A second endpoint, with a
free numerical exponent $i$, verifies the tree constructed by repeating
the selected context $i$ times. This separation prevents a hidden
unbounded universal quantifier from entering the induction invariant.
The analogous regular construction uses repeated run states.

\paragraph{Noncontracting membership.}
The endpoint replaces the arbitrary derivation existential by reachability
in the finite graph on words of length at most $n$. A shortest witness has
fewer than $V=\sum_{j\le n}m^j$ steps. For example, after numbering graph
vertices below $V$, every path with at most $V$ vertices has a list code
bounded by a primitive-recursive expression in $V$; include its length
and the step-rule witnesses in that bound. No efficient space bound is
asserted for this exhaustive algorithm.

\paragraph{Halting refutation.}
The target denies the conjunction of the displayed Boolean totality formula
and correctness as a decider for $K$. The derivation uses totality at the
explicit code of the constructed machine and a case split on its Boolean
output. It does not assert that all externally total machines have
iPRA proofs of totality. The complement refutation instead produces the
intuitionistically impossible equivalence $A\leftrightarrow\neg A$.

\subsection{Dependency order and proof obligations}

\begin{longtable}{@{}p{0.10\textwidth}p{0.40\textwidth}p{0.43\textwidth}@{}}
\toprule
Stage & Deliverable & Main proof invariant or bound\\
\midrule\endhead
1 & Words, list and table codecs & Digit identities, range tests, and
round trips on valid codes.\\
2 & Bits, least witnesses, graph paths & Bounded-formula reflection;
finite counts and path deletion.\\
3 & DFA and $\eps$-NFA runs & Equal prefix states; state-set bit equality;
empty blocks shorter than state count.\\
4 & Closure and Kleene conversions & Boundary states, finite factor lists,
expression-interval tables, elimination round index.\\
5 & Decisions and minimization & Product reachability, $n^2$ pair table,
access-word and separator certificates.\\
6 & Finite monoids & Transformation-vector equality and bounded closure
among at most $n^n$ transformations.\\
7 & Grammar trees and normalization & Frontier equality, nullable ranks,
unit paths, and local rule provenance.\\
8 & CYK and direct chart & Smaller intervals or smaller insertion ranks;
bounded witness reconstruction.\\
9 & PDA and pumping & Proper stack segments, finite endpoint chains,
longest-path repeated variables, context-splice checks.\\
10 & Noncontracting and recognizable & Length-bounded graph; supplied
trace checker; explicit lists of certificates.\\
11 & Diagonal programs & Compiler simulation and input/output traces;
Boolean branch contradictions.\\
12 & Counting & Finite bijections, matrix convolution, determinant
involutions, first-crossing reflection.\\
\bottomrule
\end{longtable}

The \texttt{lean/} directory makes these obligations concrete.
\texttt{SemanticCore.lean} provides a small candidate semantic core with
ordinary proof scripts and no placeholders, but remains uncompiled.
The subject outlines specify structures, algorithms, and detailed
remaining theorem targets. \texttt{ObjectTheoryOutline.lean.txt} lists
the fixed PR symbols and the intended closed \texttt{Prf PRAAx} endpoints.
An obligation record in an outline is a specification whose value must
still be constructed; declaring that record does not establish its fields.

\subsection{Reuse from Mathlib and from ProveIt}

Mathlib's inspected documentation supplies semantic definitions and results
for DFAs, NFAs, empty-edge automata, regular expressions, and CFG derivations
\cite{MathlibDFA,MathlibNFA,MathlibEpsilon,MathlibRegex,MathlibCFG}.
These interfaces are useful for cross-checking semantics. They often use
arbitrary types or sets and therefore need finite Boolean-table adapters
before use in this arithmetic development. At inspection, the regular
expression documentation explicitly did not supply the entire
regex--automaton equivalence; this plan therefore includes both
conversions as work to implement, rather than assuming a nonexistent
version-independent theorem name.

ProveIt's inspected arithmetic infrastructure has the stronger object-level
judgment needed here. Its bounded-formula compiler is recursion on a
\emph{fixed syntax object in the metatheory}; its output is an arithmetic
term and an arithmetic derivation. This is compatible with
Proposition~\ref{base:noeval}, because it is not a universal primitive-recursive
evaluation function in the object theory. The same distinction must be
preserved when building our new program and proof compilers.

\subsection{Completion gates and exact current status}

The inspected repository uses \texttt{leanprover/lean4:v4.32.0}. Complete
formalization should be carried out in a checkout pinned to the recorded
revision, preserving that checkout's dependency lockfile. The necessary
gates are:
\begin{enumerate}
\item Typecheck every new module with \texttt{lake env lean} or the
corresponding Lake build target.
\item Produce the explicitly named first-order formula and a value of
\texttt{Prf PRAAx []} for every claimed object-theory endpoint.
\item Inspect the endpoint formulas for validity hypotheses, empty-input
branches, nonerasing hypotheses, totality data, and certificate bounds.
\item Audit proof dependencies for \texttt{sorryAx} and added axioms.
Metatheoretic tools used to construct a derivation must not be confused
with extra axioms of the object theory.
\item If desired, transport closed results with
\texttt{IPRABridge.conciseToStandard} and inspect the translated symbols.
\end{enumerate}
No Lean executable was available in the present workspace, and none of
these Lean build gates was executed. The finite reference checks in the
next section were executed. The human proofs, unchecked Lean blueprint,
and executable checks are deliberately reported with their own statuses.
